%\documentclass{article}
%\usepackage[utf8]{inputenc}

\documentclass[12pt]{article}
\usepackage{graphicx} % This lets you include figures
\usepackage{hyperref} % This lets you make links to web locations
\graphicspath{ {./images/} }

\usepackage[rightcaption]{sidecap}
\usepackage{subcaption}
\usepackage{wrapfig}

\usepackage{float}

\usepackage{imakeidx}

\makeindex


\title{Manual}
\author{}
\date{\today}

\begin{document}
\maketitle{}

\tableofcontents

\clearpage
\newpage

\section{Flatcam:}

%\maketitle{Setup}

\subsection{calculate tool diameter}
tool→calculator
input tip-Diameter as 0.2mm
input tip angel as 60
input desired cut-z as 0.1 
use the tool diameter given in isolating routing

\subsection{Isolation Routing}
File→Open→Open Gerber
 in Project tab select the gerber file that has your isolation routing 
 select isolation routing under Tools in gerber object isolation 
 In isolation tool tool add from DB and input the tool Dia that was calculated in the tool calculator 

press the generate geometry button

in parameters for Tool 1, select the dia of the v-tip and input 0,02 then input the angle of the v-tip 60, both both x-y and z feed rates between 0-170 (170 recommended) 
set the spindle speed at 11000. 
press Generate CNCJob object 
save the cnc code for the usb drive.

\subsection{Drilling}
File→Open→Open Excellon
select drilling Tool

select the tools you want to create the CNC job fore 

Adjust Drill Z to the thickness of the copper laminate( value must be negative )

set travel Z to 2.00 mm 

set feed rate between 0 and 170(170 recommended) 
set spindle speed to 11000.

click Create CNC Job object 

save the cnc code for the usb drive 



\subsection{Board Cutout}

File→→OpenOpen Gerber for your board cutout 
select your board outline 
 Under Tools select cutout tool
specify a Margin. This will create a rectangular cutout at the given distance from any element in the Gerber. Specify a Gap Size. 2 times the diameter of the tool you will use for cutting is a good size. Specify how many and where you want the Gaps along the edge, 2 (top and bottom), 2 (left and right) or 4, one on each side.

Click on Generate Geometry.

Create a CNC job for the newly created geometry 





\subsection{2-side PCB}
open File→Open→Open Gerber Select the side that has not bin used (copper top/bottom) 

select Tools→2-sided PCB tool
under type select gerber and the side of the pcb you are working on
Then at reference select box and under reference object select gerber and board outline (depends on file name) and then select mirror.

Select drill dia (usually 1 or 1.5mm).
Select alignment drill coordinates by right click att the coordinate than press +aff 
Generate excellon object.

Go to Project select the gerber file that was mirrored proceed as in isolation routing 

Go to Project select the alignment hole file(name my differ) proceed as in drilling





\clearpage
\newpage

\section{Wegster operation:}


\subsection{start up}

load the g-code files that you wish to use onto the computer in a way that you can find them. Open Wegstr CNC controlling software if it is not already open.
if the indicator under status does not read device ready push the black button next to the power button on the wegstr CNCMilling Machine 
clamp down the copper laminate to the machine bed 
locate the  bits with corresponding dimensions  as the tools used in FlatCAM

\subsection{safety}
The manufacturer recommends the use of eye protection and protective gloves while operating the machine. 

\subsection{Isolation Routing}
change the bit to the spital bit
load the g-code for your isolation routing push the button in the software under the text auto level 
    than press the button labeled start 
use the vacuum to remove the fiberglass recommended to use while the Routing is tacking place  
\subsection{Drilling}
load the g-code for the drilling

change to the drill bit 

if the isolation routing has bin used, use the x and y movement controls to move the spindle to x = 0 and y = 0 and than use the z axis control to move the drill bit down to the copper and set z = 0 
    
else push the auto level button \
  
 than press the button labeled start
 use the vacuum to remove the fiberglass recommended to use while the drilling is tacking place  


\subsection{Board Cutout}
this is the last operation for the PCB  
change the bit to the Flat end milling bit 
load the g-code for the board cutout 
move the spindle to 0,0,0
press start 
 use the vacuum to remove the fiberglass recommended to use while the milling is tacking place  

\subsection{2-side PCB}
load the alignment drill  g-code
change the bit to your drill bit 
use the x and y movement controls to move the spindle to x = 0 and y = 0 and then use the z axis control to move the drill bit down to the copper and set z = 0 
press the start button 

declamp the laminate 
flip the laminate to the other side 
use the alignment holes to align the laminate 
clamp the laminate down to the work bed 
load the bottom isolation routing g-code 
change to the spiral bit 
move the spindle to 0,0,0 and press auto level 
press start button 
\newpage


\section{Links:}

\subsection{flatcam:}

\url{http://flatcam.org/manual/index.html}

\subsection{wegstr manual pdf:}

\url{https://testglassware.com/wp-content/uploads/2022/05/MANUAL-ITEM-1.pdf}

\newpage

%\listoffigures

\end{document}
